\documentclass[12pt,a4paper]{ctexart}
\usepackage[a4paper,margin=2.3cm]{geometry}
\usepackage{amsmath,amssymb}
\usepackage{booktabs}
\usepackage{array}
\usepackage{enumitem}
\usepackage{setspace}
\usepackage{titlesec}
\usepackage{hyperref}
\usepackage{fontspec}
\setmainfont{Times New Roman}

\setstretch{1.2}
\titleformat{\section}{\centering\bfseries\large}{}{0pt}{}
\titleformat{\subsection}{\bfseries\normalsize}{}{0pt}{}

\begin{document}

\begin{center}
    {\LARGE \textbf{Basic Principles of Marxism}}
\end{center}

\section{I. Core Definitions}

\begin{enumerate}[leftmargin=2em]
    \item \textbf{Marxism}: Marxism is a scientific theoretical system founded by Marx and Engels and continuously developed by their successors; a doctrine about the general laws governing the development of nature, society, and human thought; a doctrine about the inevitable replacement of capitalism by socialism and the eventual realization of communism; a doctrine about the liberation of the proletariat, all humanity, and the free and all-round development of every individual; and the guiding ideology of proletarian parties and socialist states as well as a guide to action for the people in creating a better life.
    
    \item \textbf{Philosophy}: Philosophy is a systematized and theorized world outlook, a summary and generalization of knowledge of nature, society, and thought, and an overall understanding of the universe and life that provides an indispensable intellectual foundation for human existence.
    
    \item \textbf{Motion}: Motion is the philosophical category that denotes the changes of all things and phenomena and their processes.
    
    \item \textbf{Contradiction}: Contradiction is the philosophical category that reflects the unity of opposites within things or between things.
    
    \item \textbf{Truth}: Truth is the philosophical category that denotes the correspondence between subjectivity and objectivity, namely the correct reflection of objective things and their laws.
    
    \item \textbf{Value}: Value is the philosophical category that reflects the meaningful relationship between subject and object.
    
    \item \textbf{Freedom}: Freedom is a category expressing the state of human activity.
    
    \item \textbf{Social formation}: Social formation is the category concerning the concrete forms, developmental stages, and qualitative states of social movement.
    
    \item \textbf{Civilization}: Civilization is the totality of the material, spiritual, and institutional achievements created by humanity, and a category marking the degree of social progress.
    
    \item \textbf{Matter}: Matter is objective reality that exists independently of human consciousness and can be reflected by human consciousness.
    
    \item \textbf{Consciousness}: Consciousness is the function and attribute of the human brain and the subjective image of the objective world.
    
    \item \textbf{Law}: A law is the inherent, internal, essential, and necessary connection in the process of the change and development of things.
    
    \item \textbf{Practice}: Practice is the social and material activity by which human beings actively transform the world.
    
    \item \textbf{Cognition}: Cognition is the active reflection of the object by the subject on the basis of practice.
    
    \item \textbf{Error}: Error is cognition that runs counter to objective things and their laws, namely a distorted reflection of objective things and their laws.
    
    \item \textbf{Social existence}: Social existence refers to the material living conditions of society and the material aspect of social life.
    
    \item \textbf{Social consciousness}: Social consciousness is the reflection of social existence and the spiritual aspect of social life.
    
    \item \textbf{Productive forces}: Productive forces are the material forces formed by human beings in production practice to transform and influence nature so that it may satisfy social needs.
    
    \item \textbf{Relations of production}: Relations of production are the economic relations formed among people in the material production process and are independent of individual will.
    
    \item \textbf{Economic base}: The economic base is the sum total of relations of production determined by productive forces at a certain stage of social development.
    
    \item \textbf{Superstructure}: The superstructure refers to ideology and the institutions, organizations, and facilities corresponding to a certain economic base.
    
    \item \textbf{Commodity}: A commodity is a product of labor produced for exchange and capable of satisfying a certain human need.
    
    \item \textbf{Commodity economy}: Commodity economy is an economic form in which production is carried out for exchange and is the product of certain social and historical conditions.
    
    \item \textbf{Law of value}: The law of value is the basic law of commodity production and commodity exchange.
    
    \item \textbf{Labor power}: Labor power is the sum total of a person's physical and mental capabilities.
    
    \item \textbf{Capital}: Capital is value that brings surplus value.
    
    \item \textbf{Surplus value}: Surplus value is the part of value created by wage laborers and appropriated without compensation by capitalists beyond the value of labor power.
    
    \item \textbf{The state}: The state is a tool by which one class rules another class.
    
    \item \textbf{The basic contradiction of capitalism}: The contradiction between the socialized nature of production and the capitalist private ownership of the means of production.
\end{enumerate}

\section{II. Basic Theoretical Explanations}

\begin{enumerate}[leftmargin=2em]
    \item \textbf{Three basic components of Marxism}: Marxist philosophy, Marxist political economy, and scientific socialism.
    
    \item \textbf{Characteristics of Marxism}: Scientific character, people's character, practical character, and developmental character.
    
    \item \textbf{The essential characteristic of Marxism}: The unity of scientific character and revolutionary character.
    
    \item \textbf{Contemporary value of Marxism}: It is a cognitive tool for observing changes in the contemporary world, a guide to action for the development of contemporary China, and scientific truth guiding the progress of human society.
    
    \item \textbf{Two major kinds of phenomena}: Material phenomena and spiritual phenomena.
    
    \item \textbf{Two major types of human activity}: Understanding the world and transforming the world.
    
    \item \textbf{The basic question of all philosophy}: The relationship between existence and thought, or the relationship between matter and consciousness.
\end{enumerate}

\section{III. Matter, Consciousness, and the Unity of the World}

\begin{enumerate}[leftmargin=2em]
    \item \textbf{The theoretical significance of the Marxist category of matter}: It insists on materialist monism and distinguishes it from idealist monism and dualism; insists on active reflection theory and the knowability of the world while criticizing agnosticism; embodies the unity of materialism and dialectics and overcomes the defects of metaphysical materialism; and embodies the unity of the materialist view of nature and the materialist view of history, thus laying the theoretical foundation for thoroughgoing materialism.
    
    \item \textbf{The fundamental attribute of matter}: The fundamental attribute of matter is motion.
    
    \item \textbf{The dialectical relationship between matter and consciousness}: Matter determines consciousness, while consciousness reacts upon matter.
    
    \item \textbf{Why artificial intelligence cannot possess human consciousness}: Human consciousness is the unity of cognition, emotion, and will, whereas artificial intelligence only simulates and extends rational intelligence; sociality is an intrinsic essential attribute of human consciousness, while artificial intelligence cannot truly possess human social attributes; human natural language is the material shell of thought and the real form of consciousness, while artificial intelligence cannot fully grasp the real meaning of natural language; and artificial intelligence can only obtain those parts of human consciousness that can be reduced to digital signals, while many aspects of the human brain cannot be reduced in this way.
    
    \item \textbf{The unity of the world}: The unity of the world lies in its materiality, and the world is unified in matter.
    \begin{enumerate}[label=(\arabic*),leftmargin=2em]
        \item Nature is material.
        \item Human society is essentially a material system formed on the basis of production practice.
        \item Human consciousness is unified with matter.
    \end{enumerate}
    
    \item \textbf{The significance of the principle of the material unity of the world}: It is the most basic and core viewpoint of dialectical materialism and the cornerstone of Marxism.
\end{enumerate}

\section{IV. Laws of Dialectics}

\begin{enumerate}[leftmargin=2em]
    \item \textbf{The general laws of development}: The law of the unity of opposites, the law of the transformation of quantity into quality, and the law of the negation of negation.
    
    \item \textbf{The law of the unity of opposites}: It is the fundamental law because it reveals the inner driving force of the development of things.
    
    \item \textbf{The laws of quantitative and qualitative change and of negation of negation}: They reveal the state, process, and trend of development.
    
    \item \textbf{The fundamental method for understanding and transforming the world}: The method of contradiction analysis.
    
    \item \textbf{The basic attributes of contradiction}: Opposition and unity.
    
    \item \textbf{The dialectical relationship between quantitative change and qualitative change}: Quantitative change is the necessary preparation for qualitative change; qualitative change is the inevitable result of quantitative change and opens the way for new quantitative change; and the two permeate each other.
    
    \item \textbf{The basic links of connection and development}: Content and form, essence and phenomenon, cause and result, necessity and contingency, and reality and possibility.
    
    \item \textbf{The nature of materialist dialectics}: Materialist dialectics is essentially critical and revolutionary.
    
    \item \textbf{The unity in dialectics}: Materialist dialectics is the unity of objective dialectics and subjective dialectics.
    
    \item \textbf{The status of contradiction analysis}: It is the methodological expression of the law of the unity of opposites, occupies the core position in the methodological system of materialist dialectics, and is the fundamental method for understanding things.
    
    \item \textbf{The living soul of Marxism}: To analyze specific problems concretely.
    
    \item \textbf{How to learn and master the scientific thinking methods of materialist dialectics}: It is necessary to strengthen dialectical thinking, historical thinking, systems thinking, strategic thinking, bottom-line thinking, and innovative thinking.
\end{enumerate}

\section{V. Practice and Cognition}

\begin{enumerate}[leftmargin=2em]
    \item \textbf{The practical character of Marxism}: Practicality is the fundamental feature that distinguishes Marxist theory from other theories; the viewpoint of practice is a basic Marxist viewpoint; and grasping the relationship between humanity and the world on the basis of practice is an important part of the Marxist world outlook.
    
    \item \textbf{The relationship between humanity and the world}: Understanding the world and transforming the world.
    
    \item \textbf{The basic characteristics of practice}: Objective reality, conscious initiative, and social-historical character.
    
    \item \textbf{The basic elements of practice}: Subject, object, and intermediary.
    
    \item \textbf{The relations between subject and object in practice}: Practical relation, cognitive relation, and value relation, among which the practical relation is the most fundamental.
    
    \item \textbf{The basic forms of practice}: Material production practice, social-political practice, and scientific-cultural practice.
    
    \item \textbf{Practice as the foundation of cognition}: Practice is the basis of cognition and plays a decisive role in cognitive activity.
    
    \item \textbf{The dialectical relationship between perceptual knowledge and rational knowledge}: Rational knowledge depends on perceptual knowledge; perceptual knowledge needs to be developed and deepened into rational knowledge; and the two permeate and contain each other.
    
    \item \textbf{The necessity and importance of returning cognition to practice}: The purpose of understanding the world is to transform it, and the truthfulness of cognition can only be tested and developed in practice.
\end{enumerate}

\section{VI. Truth and Value}

\begin{enumerate}[leftmargin=2em]
    \item \textbf{The objectivity of truth}: The objectivity of truth determines the unicity of truth.
    
    \item \textbf{The unity of the absolute and the relative in truth}: The absoluteness of truth refers to the certainty of the unity of subjectivity and objectivity and the infinity of development; the relativity of truth refers to the limited and incomplete character of human understanding of objective things and their laws under certain conditions; the two are dialectically unified; and their unity is rooted in the contradiction between the infinity and finiteness of human cognitive ability.
    
    \item \textbf{Why practice is the sole criterion for testing truth}: Truth is the correct reflection of objective things and their laws, and its nature lies in the correspondence between subjectivity and objectivity; moreover, practice has direct reality, which is the concrete manifestation of its objective reality.
    
    \item \textbf{The basic characteristics of value}: Subjectivity, objectivity, multidimensionality, and social-historical character.
    
    \item \textbf{Understanding the world}: The subject actively reflects the object, obtains scientific knowledge of the essence and development laws of things, and explores and masters truth.
    
    \item \textbf{Transforming the world}: Human beings change the existing forms of things according to the needs of their own survival and development, thus creating their ideal world and way of life.
    
    \item \textbf{The fundamental goal of understanding and transforming the world}: To understand necessity and strive for freedom.
\end{enumerate}

\section{VII. Historical Materialism}

\begin{enumerate}[leftmargin=2em]
    \item \textbf{The basic question of the view of social history}: The relationship between social existence and social consciousness.
    
    \item \textbf{Main contents of social existence}: Natural geographical environment, population factors, and the mode of material production.
    
    \item \textbf{The dialectical relationship between social existence and social consciousness}: Social existence and social consciousness are dialectically unified. Social existence determines social consciousness, while social consciousness reflects and reacts upon social existence. Social existence is the objective source of the content of social consciousness, and social consciousness is the subjective reflection of the process and conditions of social material life. Social consciousness is also the product of people's social material intercourse and is historical and concrete.
    
    \item \textbf{The fundamental division between historical materialism and historical idealism}: Whether social existence determines social consciousness or social consciousness determines social existence.
    
    \item \textbf{The basic contradictions of human society}: The contradiction between productive forces and relations of production, and the contradiction between economic base and superstructure.
    
    \item \textbf{The basic laws of the development of human society}: The law governing the movement of the contradiction between productive forces and relations of production, and the law governing the movement of the contradiction between economic base and superstructure.
    
    \item \textbf{Characteristics of productive forces}: Productive forces possess objective reality and social-historical character.
    
    \item \textbf{The basic elements of productive forces}: Means of labor, objects of labor, and labor power.
    
    \item \textbf{Science and technology}: Science and technology are the concentrated embodiment and primary mark of advanced productive forces, and science and technology are the primary productive force.
    
    \item \textbf{The relationship between productive forces and relations of production}: Productive forces determine relations of production, while relations of production react upon productive forces.
    
    \item \textbf{The political superstructure}: The political superstructure occupies the dominant position in the whole superstructure, and state power is its core.
    
    \item \textbf{The dialectical relationship between economic base and superstructure}: The economic base determines the superstructure, while the superstructure reacts upon the economic base.
    
    \item \textbf{The development and transformation of the mode of production}: It is the foundation of the formation and development of world history.
    
    \item \textbf{The social basic contradictions}: Social basic contradictions are the fundamental driving force of historical development.
    
    \item \textbf{The role of social basic contradictions in historical development}: Productive forces are the most fundamental driving force in the movement of social basic contradictions and the ultimate decisive force in social progress; social basic contradictions, especially the contradiction between productive forces and relations of production, determine the existence and development of other contradictions in society; and social basic contradictions influence and promote changes and development of social formations through different forms of manifestation and resolution.
    
    \item \textbf{The scientific and technological revolution}: It is a powerful lever promoting economic and social development.
    
    \item \textbf{The people}: The people are the subjects of social history and the creators of history.
\end{enumerate}

\section{VIII. Political Economy of Capitalism}

\begin{enumerate}[leftmargin=2em]
    \item \textbf{The two factors of commodities}: Use value and value.
    
    \item \textbf{The basic contradiction of commodity economy}: The contradiction between private labor and social labor.
    
    \item \textbf{The basic elements of the labor process producing material goods}: The labor of workers, the object of labor, and the means of labor.
    
    \item \textbf{The nature of the capitalist state}: In essence, the capitalist state is a tool for the bourgeoisie to exercise class rule.
    
    \item \textbf{The basic principles of the constitution of capitalist states}: The principle of private ownership, the principle of popular sovereignty, and the principle of separation and balance of powers.
    
    \item \textbf{The objective law of capitalist development}: Free competition leads to concentration of production and concentration of capital, and when these reach a certain stage, monopoly inevitably arises.
    
    \item \textbf{Manifestations of economic globalization}: Globalization of production, globalization of trade, and globalization of finance.
    
    \item \textbf{Limitations of capitalism}: The basic contradiction of capitalism obstructs the development of social productive forces; polarization in the possession of wealth under capitalism leads to economic crises; and the bourgeoisie dominates and controls the development and operation of the capitalist economy and politics, constantly intensifying social contradictions and conflicts.
\end{enumerate}

\section{IX. Historical Inevitability of Socialism}

\begin{enumerate}[leftmargin=2em]
    \item \textbf{The historical inevitability of socialism replacing capitalism}: The basic contradiction of capitalism contains the seeds of all modern conflicts; capital accumulation intensifies the basic contradiction of capitalism and ultimately negates capitalism itself; state-monopoly capitalism is a higher form of the socialization of capital and will become the prelude to socialism; and capitalist society contains the contradiction and struggle between the bourgeoisie and the proletariat.
    
    \item \textbf{The five-hundred-year historical process of socialism}: From utopia to science in 1848 with the \textit{Manifesto of the Communist Party}; from theory to practice in 1871 with the Paris Commune; from ideal to reality in 1917 with the October Revolution; and from one country to multiple countries in the period after the Second World War.
    
    \item \textbf{How to explore a development road suited to a country's national conditions}: It is necessary to maintain a scientific attitude toward Marxism, proceed from concrete historical conditions and follow one's own road, and fully absorb all outstanding achievements of human civilization.
    
    \item \textbf{The objective laws that must be followed in advancing socialism in practice}: The laws governing the development of human society, the laws governing socialist construction, and the laws governing Communist Party rule.
\end{enumerate}

\end{document}